\section{Retaining degreewise credit with conditional projections}\label{sec:retain-credit}
For a function $f$ of the independent letters, let $E_t f$ average over $w_t$ while fixing all other letters, and let $D_t=I-E_t$. Define
\[
 f_S=\left(\prod_{t\in S}D_t\right)\left(\prod_{t\notin S}E_t\right)f,\qquad
 (f)_r=\sum_{|S|=r}f_S.
\]
These operators commute. Expanding $\prod_t(E_t+D_t)=I$ gives $f=\sum_S f_S$. Each $f_S$ depends only on $S$ and has conditional mean zero in each of its letters. If $S\ne U$, choose a letter in their symmetric difference and average it first; this proves $\mathbb E\langle f_S,g_U\rangle=0$. Thus the decomposition is orthogonal. Its uniqueness follows by applying its defining projection to any decomposition with the same conditional-mean properties. This construction is the familiar orthogonal interaction decomposition, retained only for the degree claims.

\begin{lemma}[Matching-degree identity and exact target spectrum]\label{lem:target-contraction}
With $A_r=\mathbb E\|(G)_r\|^2$ and $J_r=\mathbb E\|(J)_r\|_F^2$,
\begin{equation}
 \mathbb E[J^\top G]=\sum_{r=0}^n\mathbb E[(J)_r^\top(G)_r],\qquad
 \|\mathbb E[(J)_r^\top(G)_r]\|\le\sqrt{A_rJ_r}.
 \label{eq:matching}
\end{equation}
The exact target masses are
\[
 \mathbb E\|G_S\|^2=\frac1K\sum_{j=1}^{K-1}|\rho_j|^{2(n-|S|)}(1-|\rho_j|^2)^{|S|},\qquad
 A_r=\binom nr\frac1K\sum_{j=1}^{K-1}|\rho_j|^{2(n-r)}(1-|\rho_j|^2)^r.
\]
In particular $\|G_S\|_{L^2}\le\lambda^{n-|S|}$.
\end{lemma}
\begin{proof}
Orthogonality eliminates all unequal supports and hence unequal degrees, giving the first identity. Cauchy--Schwarz after testing any unit parameter direction gives its bound. For the target, expand
\[
 \alpha^{js_n}=\alpha^{js_0}\prod_t\bigl[\rho_j+(\alpha^{jb_t}-\rho_j)\bigr].
\]
The term indexed by $S$ is $\alpha^{js_0}\rho_j^{n-|S|}\prod_{t\in S}(\alpha^{jb_t}-\rho_j)$. It is centered in every included letter and independent of excluded ones, so it is exactly the projection component. Its squared $L^2$ norm is $|\rho_j|^{2(n-|S|)}(1-|\rho_j|^2)^{|S|}$ by independence. Parseval across the nonconstant Fourier modes gives the vector formula, and summing the $\binom nr$ supports gives $A_r$. The convention is $0^0=1$.
\end{proof}

\begin{proposition}[Score count and degree]
\label{prop:score-degree}
Fix the start state and the query. For $\bm\theta\in\initball$, $\Jz=\sum_{k=0}^{k_0}\mathbf J^{[k]}+\mathbf J_{\mathrm{rem}}$, and there is $\attconst$ such that:
\begin{enumerate}[label=(\roman*),leftmargin=*,itemsep=2pt,topsep=2pt]
\item $\sup\|\mathbf J^{[0]}\|_F\le C\initscale$, $\ \sum_{|S|=1}\sup\|\mathbf J^{[0]}_S\|_F\le C\initscale^2$, and $\mathbf J^{[0]}_S=0$ for $|S|\ge2$;
\item for $k\ge1$, $\ \sum_S\sup\|\mathbf J^{[k]}_S\|_F\le C\initscale^2(\attconst\initscale^4)^k$, and $\mathbf J^{[k]}_S=0$ for $|S|>2k+1$;
\item $\sup\|\mathbf J_{\mathrm{rem}}\|_F\le C\initscale^2(\attconst\initscale^4)^{k_0+1}$.
\end{enumerate}
The same bounds hold for $\bm z^{[k]}$ and $\bm z_{\mathrm{rem}}$ with one more factor $\initscale$. Suprema are over inputs; constants do not depend on
$n$, $k$ or $k_0$.
\end{proposition}

\begin{proof}[Proof of \cref{prop:score-degree}]
For a term $\term$ supported on the sites $U$,
\[
 \term_S=0\quad(S\not\subseteq U),\qquad \sup\|\term_S\|\le2^{|S|}\sup\|\term\|,
\]
so
\[
 \sum_S\sup\|\term_S\|\le\sum_{S\subseteq U}2^{|S|}\sup\|\term\|=3^{|U|}\sup\|\term\| .
\]

\noindent(i) At the answer position the direct input is fixed; apply \cref{prop:network}(iii).

\smallskip\noindent(ii) For $k\ge1$, \cref{prop:network}(ii)--(iii) give $|U|\le2k+1$, hence $\mathbf J^{[k]}_S=0$ for $|S|>2k+1$, and, with $\attconst\ge9C_0$,
\begin{alignat*}{2}
 \sum_S\sup\|\mathbf J^{[k]}_S\|_F&\le\sqrt K\,3^{2k+1}\|\partial\bm z^{[k]}\|_{\mathrm{abs}} &\qquad& \why{support bound}\\
 &\le C\initscale^2(\attconst\initscale^4)^k && \why{\cref{prop:network}(ii)}
\end{alignat*}

\noindent(iii) Apply \cref{prop:network}(ii) with $k=k_0+1$. Logit bounds carry one more factor $\initscale$ than Jacobian bounds.
\end{proof}


\begin{proof}[Complete proof of \cref{thm:main}]
Set $N=\lceil(n-1)/2\rceil$ and $b=C_{\rm att}\varepsilon^4$. Choose $c$ so that $\varepsilon^2\le c\lambda$ implies $b\le\lambda^2/2$. For every support $S$,
\[
 \|\mathbb E[(J^{[k]})_S^\top G_S]\|
 \le\lambda^{n-|S|}\sup\|(J^{[k]})_S\|_F.
\]
The zero-order coefficient is additive, so its projections at degrees at least two vanish. For $k\ge1$, \cref{prop:score-degree} bounds the sum over supports by $C\varepsilon^2b^k\lambda^{n-2k-1}$. Summing over $k$ gives
\[
 C\varepsilon^6\lambda^{n-3}\sum_{k=1}^N(b/\lambda^2)^{k-1}
 \le C\varepsilon^6\lambda^{n-3}.
\]
For the remainder let $Q_{\ge2}$ denote the orthogonal projection onto degrees at least two. Then
\[
 \mathbb E[(Q_{\ge2}J_R)^\top Q_{\ge2}G]
 =\mathbb E[J_R^\top Q_{\ge2}G].
\]
Its norm is at most $\|J_R\|_{L^2(F)}\|Q_{\ge2}G\|_{L^2}\le C\varepsilon^2b^{N+1}$. As in the direct summation proof this is at most $C\varepsilon^6\lambda^{n-3}(b/\lambda^2)^N$. This proves the aggregate interaction bound. Replacing $Q_{\ge2}$ by any single degree projection gives the same bound degree by degree, with the same constant. No estimate of a remainder's supremum after projection is needed.
\end{proof}

\begin{lemma}[Gradient split]\label{lem:gradient-split}
For $p_\theta=\operatorname{softmax}(z_\theta)$,
\[
 \nabla_\theta L_{\rm out}=\mathbb E[J^\top(p_\theta-u)]-\nabla_\theta\outcredit,
 \qquad \nabla_\theta\outcredit=\mathbb E[J^\top G].
\]
\end{lemma}
\begin{proof}
Differentiate the cross-entropy, whose logit derivative is $p_\theta-y$, and split this derivative as $(p_\theta-u)-G$. Finite data support and local smoothness justify differentiation inside the expectation.
\end{proof}

\subsection{The local innovation}
For a uniform start, every $s_{t-1}$ is uniform and independent of the entire letter sequence. Let $q_s=(1-p)e_s+pe_{s+1}$ and
\[
 g_t=e_{s_t}-q_{s_{t-1}}=(b_t-p)(e_{s_{t-1}+1}-e_{s_{t-1}}).
\]
Its mean is zero. In a unitary Fourier basis on $\bm1^\perp$, its covariance has eigenvalues
\[
 \frac{p(1-p)}K|\alpha^j-1|^2
 =\frac{1-|\rho_j|^2}K\ge\frac{1-\lambda^2}K.
\]
Indeed the uniform state gives covariance $I/K$ for its Fourier components, and the independent centered match bit contributes $p(1-p)$. This proves the local innovation inequality without a degree decomposition. Define $\rulecredit=\mathbb E_t\mathbb E\langle z_t,g_t\rangle$.

\subsection{Step rule credit}

Let $\bm g_t=(b_t-p)(\bm e_{s_{t-1}+1}-\bm e_{s_{t-1}})$ and $\rulecredit(\bm\theta)=\E_t\E[\langle\bm z_t,\bm g_t\rangle]$, $t$ uniform on $\{2,\ldots,n\}$. Letter $w_t$ sits at
$l_t=1+t$, and $i_t=n+1+t$ holds $s_{t-1}$ and predicts $s_t$, so $m_{l_t}=t+2$ and $m_{i_t}=n+2+t$. Write $\kappa_t=\kappa_{l_t,i_t}$. The displayed state tokens visible from $i_t$ are functions of $(s_0,w^\star,w_{<t})$; the prompt also includes the letters. The current letter is therefore the only visible source that depends on $b_t$ after conditioning on $s_0,w^\star,w_{\ne t}$.

\begin{lemma}[Average dilution]
\label{lem:dilution-average}
For $n\ge2$,
\begin{equation}
 \E_t[\kappa_t]\le\frac{1}{n+4}\Bigl(1+\log2+\frac{n+2}{n-1}\Bigr)\le\frac{6}{n} .
 \label{eq:dilution-avg}
\end{equation}
\end{lemma}

\begin{proof}
$m_{i_t}\ge n+4$ and $\sum_{l_t<j\le i_t}m_j^{-1}\le\log2+\log\frac{n+2}{t+2}$; also $\sum_{t=2}^n\log\frac{n+2}{t+2}\le\sum_{k=1}^{n+2}\log\frac{n+2}{k}\le n+2$.
\end{proof}

\begin{lemma}[Rule-credit sites]
\label{lem:rule-sites}
Under (A3$'$), a term $\term$ of $\bm z_t$ or $\bm h^2_{i_t}$ has $\E\langle\term,\bm g_t\rangle=0$ unless it contains both the letter input $\bm x_{l_t}$ and the input of a
state ($s_0$ or a displayed state). In particular $\bm z^{[0]}_t$ contributes nothing.
\end{lemma}

\begin{proof}
If $\bm x_{l_t}\notin\term$, then $\E[\bm g_t\mid s_0,w^\star,w_{\ne t}]=0$. If $\term$ contains no state input, it is a function of $(w^\star,w)$ and, by (A3$'$),
$\E[\bm g_t\mid w^\star,w]=(b_t-p)(\bm u-\bm u)=0$. A contributing term has two sites, so $k\ge1$ by \cref{prop:network}(iii).
\end{proof}

For readout indices $a,b$ let
\begin{gather*}
 Q_t=\sum_{\ell=1,2}\sum_{\mathrm{heads}}\frac{1}{m_i}\sum_{j\le i}\mathbf W_V\bm x_j\Bigl(\hat\psi_{ij}-\frac{1}{m_i}\sum_{l\le i}\hat\psi_{il}\Bigr),\qquad i=i_t,\\
 \hat\psi_{ij}=\frac{(\mathbf W_Q\bm x_i)^\top\mathbf W_K\bm x_j}{\sqrt{d_{\mathrm h}}},
\end{gather*}
and $P_n(\tilde{\bm\theta})=\E_t\E[(\bm g_t)_a(Q_t)_b]|_{\bm\theta=\tilde{\bm\theta}}$.

\begin{lemma}[Leading rule coefficient]
\label{lem:leading}
Assume (A1) with full-prefix attention, (A2) and (A3$'$). Fix a readout entry $V_{ab}$ and let
\[
 F_n(\initscale)=\pd{\rulecredit}{V_{ab}}\Big|_{\bm\theta=\initscale\tilde{\bm\theta}}=\E_t\E[(\bm g_t)_ah^2_{i_t,b}]\big|_{\bm\theta=\initscale\tilde{\bm\theta}},\qquad\text{so }\ \|\nabla_{\bm\theta}\rulecredit(\bm\theta_0)\|\ge|F_n(\initscale)|.
\]
\begin{enumerate}[label=(\roman*),leftmargin=*,itemsep=0pt,topsep=2pt]
\item $P_n$ is a homogeneous polynomial of degree $6$ in the attention and embedding entries of $\tilde{\bm\theta}$, independent of the activation
pattern, and for $\initscale\le\initscale_0$ with $\attconst\initscale^4\le\frac{1}{2}$,
\[
 F_n(\initscale)=\initscale^6P_n(\tilde{\bm\theta})+r_n(\initscale),\qquad |r_n(\initscale)|\le\frac{C}{n}\Bigl(\frac{\initscale^7}{\initscale_0}+\initscale^{10}\Bigr).
\]
\item For a fixed query with $0<p<\frac{1}{2}$, $P_n$ contains a degree-$6$ count--letter monomial with coefficient at least
$\dfrac{p(1-p)\min\{1-2p,\frac{1}{4}\}}{3Kn\sqrt{d_{\mathrm h}}}>0$ for $n\ge2$; in particular $P_n\not\equiv0$.
\end{enumerate}
\end{lemma}

\begin{proof}
(i) Let $k_0=2+\lceil\log_2n\rceil$. By \cref{lem:rule-sites}, the zero-score term contributes nothing; by \cref{prop:network} and
\cref{lem:dilution-average}, the terms with at least two scores and the remainder contribute at most
\[
 C\initscale^2\E_t[\kappa_t]\sum_{k=2}^{k_0}(\attconst\initscale^4)^k+C\initscale^2(\attconst\initscale^4)^{k_0+1}\le\frac{C\initscale^{10}}{n} .
\]
With the activation matrices fixed, a one-score monomial satisfies $\term(\initscale\tilde{\bm\theta})=\initscale^{\deg\term}\term(\tilde{\bm\theta})$ with $\deg\term\ge6$, and the
degree-$6$ part is exactly $Q_t$. An MLP matrix correction adds two degrees and a bias output in place of an input adds one, so the
remaining one-score terms are $O(\initscale^7)$; with $\initscale_0\le(2N_b)^{-1}$ and \cref{lem:dilution-average}, their absolute sum is at most
$\frac{C\initscale^7}{n\initscale_0}$, with $C$ independent of $N_b$.

(ii) Choose a block-$2$ head, a state $s^\ast$, an output coordinate $b$, and $a=s^\ast+1$. Set
$\tilde E_{s^\ast,0}=\tilde E_{w^\star,1}=\tilde E_{w^\star,2}=(\tilde{\mathbf W}_Q)_{0,0}=(\tilde{\mathbf W}_K)_{0,1}=(\tilde{\mathbf W}_V)_{b,2}=1$ and all other entries to $0$; the value of $P_n$ is then the
coefficient of the product of these six entries. With $N_t$ the number of visible query tokens,
\begin{gather*}
 (Q_t)_b=\frac{\mathbf 1[s_{t-1}=s^\ast]}{\sqrt{d_{\mathrm h}}}\Bigl(\frac{N_t}{n+2+t}-\frac{N_t^2}{(n+2+t)^2}\Bigr),\\
 \E[(b_t-p)N_t]=p(1-p),\qquad \E[(b_t-p)N_t^2]=p(1-p)\bigl[3+2(n-1)p\bigr].
\end{gather*}
Under (A3$'$), $s_{t-1}$ is independent of the letters, so the coefficient is
\[
 \frac{1}{K\sqrt{d_{\mathrm h}}(n-1)}\sum_{t=2}^n\frac{p(1-p)}{n+2+t}\Bigl(1-\frac{3+2(n-1)p}{n+2+t}\Bigr)
 \ge\frac{p(1-p)\min\{1-2p,\frac{1}{4}\}}{3Kn\sqrt{d_{\mathrm h}}},
\]
using $3+2(n-1)p\le\max\{2p,\frac{3}{4}\}(n+4)$ for $p<\frac{1}{2}$, $n\ge2$, and $n+2+t\le3n$.
\end{proof}

\begin{proof}[Proof of \cref{thm:separation}(i)]
\emph{Step gradient.} At the position $i$ predicting $s_t$, $\bm h^{(0)}_i=\mathbf E_{\mathrm{val}}^\top\bm e_{s_{t-1}}+\mathbf P\bm e_i$. Let
$\bm\Gamma=\frac{1}{K}\sum_{s}\bigl((1-p)\bm e_s+p\bm e_{s+1}-\bm u\bigr)\bm e_s^\top$. For $\initscale\le\frac{\mixrate\embconst}{4KC}$,
\begin{alignat*}{2}
 \bm h_i&=\bm h^{(0)}_i+\bm O^2+O(\initscale^2),\qquad \bm p_i-\bm u=O(\initscale^2) &\qquad& \why{\cref{prop:network}(i)}\\
 \pd{\Lstep}{\mathbf V}&=-\E\bigl[(\bm e_{s_t}-\bm u)\bm h^{(0)\top}_i\bigr]+O(\initscale^2) && \why{$\bm O^2$ fixed, $\E[\bm e_{s_t}-\bm u]=0$}\\
 &=-\bm\Gamma\mathbf E_{\mathrm{val}}+O(\initscale^2) && \why{uniform start}\\
 \|\nabla_{\bm\theta}\Lstep\|&\ge\|\bm\Gamma\mathbf E_{\mathrm{val}}\|_F-C\initscale^2 && \\
 &\ge\frac{\mixrate}{K}\cdot\frac{\initscale\embconst}{2}-C\initscale^2\ge\frac{\initscale\mixrate\embconst}{4K} && \why{$\sigma_{\max}(\bm\Gamma)=\tfrac{\mixrate}{K}$, $\sigma_{\min}(\mathbf E_{\mathrm{val}})\ge\tfrac{\initscale\embconst}{2}$}
\end{alignat*}
using that $\mathbf E_{\mathrm{val}}$ has full row rank.

\emph{Outcome gradient.} On $\initball$ and for $\initscale\le c\mixrate$,
\begin{alignat}{2}
 \nabla_{\bm\theta}\Lout&=\E\bigl[(\Jz)^\top(\bm p-\bm u)\bigr]-\nabla_{\bm\theta}\outcredit &\qquad& \why{\cref{lem:gradient-split}}\nonumber\\
 \|\nabla_{\bm\theta}\Lout\|&\le C\initscale\cdot C\initscale^2+C\initscale\mixrate^n\le C\initscale(\initscale^2+\mixrate^n) && \why{\cref{prop:network}(i), \cref{lem:score-sum}}\label{eq:outcome-grad}
\end{alignat}
\end{proof}

\begin{proof}[Proof of \cref{thm:separation}(ii)]
Let $k_0=1+\lceil\log_2n\rceil$ and $\attconst\initscale^4\le\frac{1}{2}$. By \cref{lem:rule-sites}, only terms with at least one score contribute.
\begin{alignat*}{2}
 \|\nabla_{\bm\theta}\rulecredit(\bm\theta)\|
 &\le2C\initscale^2\,\E_t[\kappa_t]\sum_{k=1}^{k_0}(\attconst\initscale^4)^k+2C\initscale^2(\attconst\initscale^4)^{k_0+1} &\qquad& \why{\cref{prop:network}(ii),(iv)}\\
 &\le C\attconst\initscale^6\bigl(\E_t[\kappa_t]+2^{-k_0}\bigr) && \why{geometric sum}\\
 &\le C\attconst\initscale^6\Bigl(\frac{6}{n}+\frac{1}{2n}\Bigr)\le\frac{C\initscale^6}{n} && \why{\cref{lem:dilution-average}}
\end{alignat*}
\end{proof}

\paragraph{Zero readout.}\label{rem:zero-readout} \Cref{thm:separation} holds with a zero readout: part~(i) uses $\bm p-\bm u=O(\initscale^2)$, exact at $\bm\theta_0$, and part~(ii) concerns $P_n$ of
\cref{lem:leading}, which contains no readout entry.


\begin{remark}[Generic initialization]
\label{rem:generic-rule}
For $p<\frac{1}{2}$, $P_n$ is a nonzero polynomial (\cref{lem:leading}(ii)), so $P_n(\tilde{\bm\theta})\ne0$ almost surely under any absolutely continuous law on
the attention and embedding entries. If $\tilde{\bm\theta}$ satisfies (A2) and $|P_n(\tilde{\bm\theta})|\ge\frac{\rulemargin}{n}$, then for
$\initscale\le\min\{\initscale_0,(2\attconst)^{-1/4},\frac{\rulemargin\initscale_0}{4C},(\frac{\rulemargin}{4C})^{1/4}\}$, with $C$ from \cref{lem:leading}(i),
\[
 \|\nabla_{\bm\theta}\rulecredit(\initscale\tilde{\bm\theta})\|\ge|F_n(\initscale)|\ge\initscale^6\frac{\rulemargin}{n}-\frac{C}{n}\Bigl(\frac{\initscale^7}{\initscale_0}+\initscale^{10}\Bigr)\ge\frac{\initscale^6\rulemargin}{2n},
\]
which matches \cref{thm:separation}(ii) up to constants. Both the threshold and the constant depend on $\rulemargin$; a lower bound on $\rulemargin$
at random initialization would need anti-concentration of $nP_n$, which we do not prove. The manuscript experiments measure the normalized credit.
\end{remark}

